(1) 固定离子正背景密度 \(n>0\)，电子电荷 \(-e\) 且 \(e>0\)，忽略电子惯性、电阻、电子压强以及位移电流。均匀背景磁场 \(\boldsymbol B_0=B_0\hat z\)，\(B_0>0\)。电子速度 \(\boldsymbol v_e=-\boldsymbol J/(ne)\)，Ohm关系 \(\boldsymbol E+\boldsymbol v_e\times\boldsymbol B=0\)，Ampère近似 \(\boldsymbol J=\nabla\times\boldsymbol B/\mu_0\)。对只依赖 \(z,t\) 的小横向磁扰动 \(\boldsymbol b\perp\hat z\)，推导线性Hall感应方程及 \(D_H\) 的量纲。
(2) 取 \(\boldsymbol b\propto(\hat x\pm i\hat y)e^{i(kz-\omega t)}\) 求两螺旋本征频率，明确正频率向 \(+z\) 传播时的极化与相速、群速；解释负频率的另一支不代表能量不稳定增长。
(3) 不依赖模式色散，以线性Hall方程推导局域正定扰动磁能 \(e_b=|\boldsymbol b|^2/(2\mu_0)\) 的连续性方程及通量 \(S_z\)。对一束严格圆偏振正频平面波，验证平均通量等于 \(v_ge_b\)，并说明Hall力为何不直接产生焦耳损耗。
(4) 用 \(d_e=\sqrt{m_e/(\mu_0ne^2)}\)、\(\Omega_e=eB_0/m_e\) 验证 \(D_H=\Omega_ed_e^2\)。讨论 \(k d_e\ll1\) 及电子惯性或有限电离子响应导致模型失效的条件，核对 \(B_0\to0\) 的形式极限。
